\section{模型假设与符号说明}

\subsection{模型假设}

为保证模型可计算且与题意一致，本文作出如下假设。

\begin{enumerate}
    \item 五年内小区空间结构、道路连通条件和服务站可选位置保持不变，小区间距离矩阵可代表老人接受服务的相对空间成本。
    \item 新增 60 岁老人进入自理状态，不直接进入半失能或失能状态；该假设符合年龄刚达到老年界限时健康状态通常较好的现实特征。
    \item 同一健康状态老人在服务需求频次上具有同质性，即同类老人采用相同月均服务需求频次。
    \item 消费能力不足时，老人按各项服务理论需求比例等比例压缩服务次数，不改变服务偏好结构。
    \item 服务站总服务能力可按日服务人次统一折算，六类服务共享站点基础服务能力。
    \item 问题二中的价格满意度按基准价格计算；问题三再根据优化定价重新计算价格满意度。
    \item 紧急救助为公益免费项目，不参与定价优化，不计政府补贴，但计入直接服务支出。
\end{enumerate}

\subsection{符号说明}

\begin{table}[H]
\centering
\caption{主要符号说明}
\label{tab:symbols}
\begin{tabular}{c p{0.68\textwidth} c}
\toprule
符号 & 含义 & 单位 \\
\midrule
$i,j$ & 小区或候选服务站位置编号，$i,j=1,2,\ldots,10$ & -- \\
$k$ & 老人健康状态编号，$k=1,2,3$ 分别表示自理、半失能和失能 & -- \\
$m$ & 服务项目编号，$m=1,2,\ldots,6$ & -- \\
$N_{i,k}^{(t)}$ & 第 $t$ 年末小区 $i$ 中第 $k$ 类老人数量 & 人 \\
$p_i,q_i$ & 小区 $i$ 中自理转半失能、半失能转失能的概率 & -- \\
$d,g$ & 老人自然死亡率和新增老人比例 & -- \\
$D_{i,m}^{\mathrm{th}}$ & 小区 $i$ 对服务 $m$ 的理论月需求量 & 次/月 \\
$D_{i,m}^{\mathrm{bd}}$ & 小区 $i$ 对服务 $m$ 的消费约束后月需求量 & 次/月 \\
$x_{j,s}$ & 是否在小区 $j$ 建设规模 $s$ 的服务站 & 0/1 \\
$z_{i,j}$ & 小区 $i$ 是否由服务站 $j$ 作为主服务站 & 0/1 \\
$w_{i,j}$ & 小区 $i$ 的服务需求由站点 $j$ 实际承接的比例 & -- \\
$Q_s$ & 规模 $s$ 的服务站日服务能力 & 人次/日 \\
$F_s,O_s$ & 规模 $s$ 的服务站建设成本和年度运营成本 & 万元、元 \\
$S_{i,j}$ & 小区 $i$ 对服务站 $j$ 的综合满意度 & -- \\
$C$ & 空间覆盖率 & -- \\
$\bar{S}$ & 老人规模加权平均满意度 & -- \\
$\rho_j$ & 服务站 $j$ 的利润率 & -- \\
$Reg_l$ & 情景 $l$ 下基准方案的后悔值 & -- \\
\bottomrule
\end{tabular}
\end{table}
